\documentclass[11pt]{article}
\usepackage{amsmath,amssymb,amsthm}
\usepackage[margin=1in]{geometry}

\newtheorem{theorem}{Theorem}
\newtheorem{lemma}[theorem]{Lemma}

\title{Problem 651: Patterned Cylinders}
\author{Project Euler Solutions}
\date{}

\begin{document}
\maketitle

\section{Problem Statement}

A cylinder is formed by a ring of $n$ cells, each coloured with one of $k$ colours. Two colourings are considered identical if one can be obtained from the other by rotation of the cylinder. Determine the number $N(n,k)$ of distinct colourings.

\section{Mathematical Foundation}

\begin{theorem}[Burnside's Lemma]
Let a finite group $G$ act on a finite set $X$. The number of distinct orbits is
\[
|X/G| = \frac{1}{|G|}\sum_{g \in G} |X^g|,
\]
where $X^g = \{x \in X : g \cdot x = x\}$.
\end{theorem}

\begin{proof}
Define the set $S = \{(g,x) \in G \times X : g \cdot x = x\}$. Counting by group elements gives $|S| = \sum_{g \in G}|X^g|$. Counting by set elements gives $|S| = \sum_{x \in X}|\mathrm{Stab}(x)|$. By the orbit-stabiliser theorem, $|\mathrm{Stab}(x)| = |G|/|\mathrm{Orb}(x)|$, so
\[
\sum_{x \in X} \frac{|G|}{|\mathrm{Orb}(x)|} = |G| \sum_{O \in X/G} \sum_{x \in O} \frac{1}{|O|} = |G| \cdot |X/G|.
\]
Dividing by $|G|$ yields the result.
\end{proof}

\begin{lemma}
Let $G = \mathbb{Z}/n\mathbb{Z}$ act on $X = \{1,\ldots,k\}^n$ by cyclic rotation. Rotation by $d$ positions fixes exactly $k^{\gcd(d,n)}$ colourings.
\end{lemma}

\begin{proof}
A colouring $(c_0, \ldots, c_{n-1})$ is fixed by rotation by $d$ iff $c_i = c_{i+d \bmod n}$ for all $i$. This forces $c_i = c_j$ whenever $i \equiv j \pmod{\gcd(d,n)}$. There are $\gcd(d,n)$ equivalence classes, each freely colourable, giving $k^{\gcd(d,n)}$ fixed colourings.
\end{proof}

\begin{theorem}[Necklace Formula]
The number of distinct necklaces is
\[
N(n,k) = \frac{1}{n}\sum_{d \mid n} \varphi(n/d)\, k^d,
\]
where $\varphi$ is Euler's totient function.
\end{theorem}

\begin{proof}
By Burnside's Lemma and Lemma~2:
\[
N(n,k) = \frac{1}{n}\sum_{d=0}^{n-1} k^{\gcd(d,n)}.
\]
Grouping by $e = \gcd(d,n)$: for each divisor $e \mid n$, there are $\varphi(n/e)$ integers $d \in \{0,\ldots,n-1\}$ with $\gcd(d,n)=e$. Hence
\[
N(n,k) = \frac{1}{n}\sum_{e \mid n} \varphi(n/e)\, k^e. \qedhere
\]
\end{proof}

\section{Algorithm}

\begin{verbatim}
function NecklaceCount(n, k, mod):
    divisors <- FindDivisors(n)
    result <- 0
    for each d in divisors:
        phi_val <- EulerTotient(n / d)
        result <- result + phi_val * ModPow(k, d, mod)
    result <- result * ModInverse(n, mod) mod mod
    return result
\end{verbatim}

\section{Complexity Analysis}

\begin{itemize}
\item \textbf{Time:} $O(\sqrt{n} \cdot \log k)$. Finding divisors costs $O(\sqrt{n})$; each of the $O(\sqrt{n})$ divisors requires $O(\sqrt{n})$ for $\varphi$ and $O(\log k)$ for modular exponentiation.
\item \textbf{Space:} $O(\sqrt{n})$ for storing divisors.
\end{itemize}

\section{Answer}

\[
\boxed{448233151}
\]

\end{document}
